\documentclass[11pt]{article}
\usepackage{fontspec}
\usepackage[french]{babel}
\usepackage{amsmath}
\usepackage{amssymb}
\usepackage{amsthm}
\usepackage{mathrsfs}
\usepackage{tikz-cd}
\usepackage{enumitem}

\newcommand{\printedpage}[1]{\marginpar{\footnotesize p.~#1}}
\newcommand{\msp}[1]{\marginpar{\footnotesize [#1]}}
\newcommand{\Loc}{\mathrm{Loc}}

\makeatletter
\renewcommand\@makefnmark{\hbox{\textsuperscript{(\@thefnmark)}}}
\renewcommand\@makefntext[1]{\noindent\textsuperscript{\@thefnmark}\,#1}
\makeatother

\begin{document}
\setcounter{footnote}{2}

\printedpage{4}
% head: A. Grothendieck : Dérivateurs, Ch. XVI, Édition des Universités de Montpellier II et Paris VII
où, on l'a vu dans $(a)$, $i_s^X$ et $i_s^Y$ sont dans $W$. Par hypothèse on a $f_s \in W$, donc $s\backslash f \in W$, q.e.d.

\textit{b}) Loc (1, $2_0$, $3'$) $\Longrightarrow$ Loc (1, 2, 3). Prouvons d'abord Loc (2). Appliquant Loc ($3'$ \textit{a}) à la projection
\[
\boldsymbol{\Delta}^1 \times X \longrightarrow X \qquad \text{(où $X \in \operatorname{Ob}\mathrm{Cat}$),}
\]
qui est fibrée scindée à fibres isomorphes à $\boldsymbol{\Delta}^1$, donc $W$-asphériques par Loc ($2_0$), on voit que cette projection est dans $W$. Cela permet d'appliquer les arguments standard d'homotopie (\textit{cf.} XII, p. 1 \ldots) pour conclure qu'un homotopisme est dans $W$, donc une catégorie contractile est $W$-asphérique, a fortiori une catégorie avec objet final (ou avec objet initial) est $W$-asphérique, \textit{i.e.} on a Loc (2).

Prouvons Loc (3). Comme Loc (1, 2, 3) $\Longleftrightarrow$ Loc (1, $2^\circ$, $3^\circ$), il suffit de prouver Loc ($3^\circ$), vu qu'on a prouvé Loc ($2^\circ$). Soit
\[
f : X \longrightarrow S \qquad \text{$W$-\textit{co}asphérique,}
\]
prouvons $f \in W$. Introduisons la catégorie \textit{fibrée} sur $S$ dont les fibres sont les $s\backslash X$, soit $\Phi_S(X)$, et considérons l'inclusion canonique
\[
i_X : X \longrightarrow \Phi_S(X).
\]
\msp{page 7}Par Loc ($3'$ \textit{b}) elle est dans $W$. D'autre part, les fibres $s\backslash X$ de $\Phi_S(X)$ sont $W$-coasphériques par hypothèse, donc par Loc ($3'$ \textit{a}) $\Phi_S(X) \longrightarrow S$ est dans $W$. Comme $i_X \in W$, il s'ensuit $(f : X \longrightarrow S) \in W$, q.e.d.

\textit{b} bis) Loc (1, $2_0$, $3'$ bis) $\Longrightarrow$ Loc (1, 2, 3 bis). On a déjà prouvé Loc (1, 2, 3) (vu que Loc ($3'$ bis) contient Loc ($3'$) comme cas particulier). Prouvons donc Loc (3 bis) en admettant déjà tout ça. Il suffira encore de le prouver sous forme duale, puisque
\[
\text{Loc (1, 2, 3 bis)} \iff \text{Loc(1, $2^\circ$, 3 bis$^\circ$),}
\]
\msp{page 8}et que Loc ($2^\circ$) est connu par Loc (1, 2, 3) $\Longleftrightarrow$ Loc (1, $2^\circ$, $3^\circ$). Il reste donc à prouver que si
\[
f : X \longrightarrow Y \qquad \text{dans } \mathrm{Cat}/S
\]
est telle que
\[
s\backslash f : s\backslash X \longrightarrow s\backslash Y
\]
est dans $W$ pour tout $s \in \operatorname{Ob} S$, alors $f \in W$. Considérons alors le carré commutatif
\[
\begin{tikzcd}
X \arrow[r, "i_X"] \arrow[d, "f"'] & \Phi_S(X) \arrow[d, "\Phi_S(f)"] \\
Y \arrow[r, "i_Y"] & \Phi_S(Y),
\end{tikzcd}
\]
où $i_X, i_Y \in W$ par Loc ($3'$ bis \textit{b}). L'hypothèse sur $f$ signifie que $\Phi_S(f)$ induit des $W$-équivalences sur les fibres, donc par Loc ($3'$ bis) $\Phi_S(f) \in W$, donc $f \in W$ par Loc (1), q.e.d.
% foot: Version 18 mars 2019

\printedpage{5}
% head: A. Grothendieck : Dérivateurs, Ch. XVI, Édition des Universités de Montpellier II et Paris VII
\section*{2\quad Axiomes Loc (4) et Loc (5) (des sommes directes, de connexité)\protect\footnote{Il s'agit ici d'axiomes qu'on peut considérer comme étant de nature plus ou moins “triviale”. Dans tous les cas à ma connaissance, ils sont bel et bien trivialement vérifiés.}}

\begin{description}[leftmargin=7em, labelwidth=6em]
\item[Loc (4)] (Axiome des \textit{sommes directes}). Pour une famille $(f_i)_{i\in I}$ de flèches dans Cat, si $f = \coprod f_i$, on a $f \in W$ si et seulement si $f_i \in W$ pour tout $i \in I$.
\end{description}

Cet axiome se décompose en deux, en le “si” et le “seulement si”, que je vais formuler sous la forme la plus économique :

\textbf{Proposition 2.1.} \textit{L'axiome} Loc (4) \textit{des sommes équivaut à la conjonction des deux axiomes suivants :}

\begin{description}[leftmargin=7em, labelwidth=6em]
\item[Loc (4 a)] \msp{page 9}Soit $(f_i)_{i\in I}$ une famille de flèches dans Cat, si les $f_i \in W$, alors $f = \coprod f_i \in W$.
\item[Loc (4 b)] Soit $S$ dans Cat, et $S_0 \hookrightarrow S$ un sommand direct (sous-catégorie à la fois ouverte et fermée). Soit $X$ dans Cat$/S$. Si $X \longrightarrow S$ est dans $W$, alors $X_0 = X|S_0 \longrightarrow S_0$ aussi. (En d'autres termes, l'inclusion $S_0 \hookrightarrow S$ est une “$W$-fibration”, \textit{i.e.} est dans $\mathrm{Fib}_W$.)
\end{description}

\textbf{Corollaire 2.2.} \textit{Considérons la variante de} Loc (4 a) \textit{obtenue en se bornant aux familles finies.} (\textit{Elle équivaut à} Loc (4) \textit{lorsque $W$ est stable par} $\varinjlim$ \textit{filtrantes,} cf. \textit{axiome} Loc (8) \textit{plus bas}). \textit{Cette variante équivaut à}

\begin{description}[leftmargin=7.5em, labelwidth=6.5em]
\item[Loc (4 a$_0$)] Tout objet $X$ de Cat est “$W$-cofibrant”, \textit{i.e.} $\emptyset \longrightarrow X$ est dans $\mathrm{Cof}_W$, \textit{i.e.} pour toute $f : Y \longrightarrow Y'$ qui est dans $W$, le morphisme correspondant $\mathrm{id}_X \amalg f : X \amalg Y \longrightarrow X \amalg Y'$ est dans $W$.
\end{description}

\textbf{Rappel 2.3.} De toutes façons, si $W$ est un localiseur fondamental (ce que nous supposerons désormais, sauf mention du contraire), l'analogue \textit{fini} de Loc (4 a) pour les produits est vrai, \textit{i.e.} si on a une famille \textit{finie} de flèches $(f_i)_{i\in I}$ dans Cat, si les $f_i$ sont dans $W$, alors $f = \prod f_i \in W$. \texttt{[Pour cela, il faut supposer que $W$ satisfait à l'axiome de localisation, forme \textit{forte}, Loc 3 bis.]} L'analogue pour une famille infinie n'est vérifié pratiquement pour aucun localiseur fondamental.

\msp{page 10}\textbf{Proposition 2.4.} \textit{Soit $W$ un localiseur fondamental. Alors les produits finis existent dans} $\mathrm{Hot}_W$\textit{, et le foncteur} $\mathrm{hot}_W : \mathrm{Cat} \longrightarrow \mathrm{Hot}_W \overset{\textit{déf}}{=} (\mathrm{Cat})W^{-1}$ \textit{y commute. Si $W$ satisfait à} Loc (4 a) (resp. Loc (4 a$_0$)), \textit{alors les sommes quelconques} (resp. \textit{finies}) \textit{existent dans} $\mathrm{Hot}_W$\textit{, et le foncteur} $\mathrm{hot}_W$ \textit{y commute. De plus, dans ce cas les sommes sont distributives par rapport aux produits (vu qu'il en est ainsi dans} Cat).

La démonstration résulte simplement du résultat général :

\textbf{Lemme 2.5.} \textit{Soit $(\mathcal{M}, W)$ une catégorie de modèles ($W$ partie quelconque de} $\mathrm{Fl}\,\mathcal{M}$). \textit{Supposons que les produits} (resp. \textit{les produits finis,} resp. \textit{les produits binaires}) \textit{existent dans $\mathcal{M}$, et que de plus pour toute famille $(f_i)_{i\in I}$ de flèches de $\mathcal{M}$} (resp. \textit{toute famille finie,}
% foot: Version 18 mars 2019
\printedpage{6}
% head: A. Grothendieck : Dérivateurs, Ch. XVI, Édition des Universités de Montpellier II et Paris VII
resp. \textit{famille binaire}) \textit{$f_i \in W$ pour tout $i \in I$ implique $f = \prod f_i \in W$. Alors les produits du type envisagé existent dans} $\mathrm{Ho}_W = \mathcal{M}W^{-1}$\textit{, et le foncteur} $\mathrm{ho}_W : \mathcal{M} \longrightarrow \mathcal{M}W^{-1}$ \textit{y commute. Enfin, si $\mathcal{M}$ admet un objet final $e$, alors $e$ est aussi objet final de} $\mathrm{Ho}_W$.

\texttt{[Il n'est pas clair que l'assertion soit vraie dans cette généralité pour les produits \textit{infinis} ; elle est vraie si $W$ est la classe des équivalences faibles d'une \textit{catégorie de modèles de Quillen.}]}

\msp{page 11}On en conclut les énoncés duaux concernant les sommes directes et les objets initiaux.

La démonstration est un cas particulier du résultat bien utile suivant, appliqué au couple de foncteurs adjoints
\[
\begin{tikzcd}
\mathcal{M} \arrow[r, shift left, "\mathrm{diag}"] & \mathcal{M}^I \arrow[l, shift left, "\text{\scriptsize\shortstack{produits\\ dans $\mathcal{M}$\\ indexés par $I$}}"]
\end{tikzcd}
\qquad\qquad :
\]

\textbf{Lemme 2.6}\footnote{\textit{cf.} XII, lemme page 248.}. \textit{Soient $(\mathcal{M}, W)$, $(\mathcal{M}', W')$ deux catégories de modèles, et}
\[
\begin{tikzcd}
\mathcal{M} \arrow[r, shift left, "\varphi"] & \mathcal{M}' \arrow[l, shift left, "\psi"]
\end{tikzcd}
\]
\textit{un couple de foncteurs adjoints compatibles aux localiseurs,} i.e. \textit{tels que}
\[
\varphi(W) \subset W', \qquad \psi(W') \subset W.
\]
\textit{Alors les foncteurs localisés}
\[
\begin{tikzcd}
\mathcal{M}W^{-1} \arrow[r, shift left, "\bar{\varphi}"] & \mathcal{M}'W'^{-1} \arrow[l, shift left, "\bar{\psi}"]
\end{tikzcd}
\]
\textit{sont également adjoints, les flèches d'adjonction pour $\bar{\varphi}\bar{\psi}$ et $\bar{\psi}\bar{\varphi}$ étant déduites de celles pour $\varphi\psi$ et $\psi\varphi$ par composition avec les foncteurs de localisation $\mathcal{M} \longrightarrow \mathcal{M}W^{-1}$, $\mathcal{M}' \longrightarrow \mathcal{M}'W'^{-1}$.}

\begin{description}[leftmargin=7em, labelwidth=6em]
\item[Loc (5)] \msp{page 12}(Axiome de \textit{connexité} ou de \textit{non-trivialité}).

Soit $f \in W$, alors $\pi_0(f)$ est bijectif.
\end{description}

Cette propriété de $W$ est d'une utilisation constante en pratique, quand on travaille avec $W$ au lieu de se borner à construire des $W$ et faire joujou avec les axiomes et leur stabilité par passage à des $\underline{\mathrm{Hom}}(I, M)$ (où ici $\mathcal{M} = \mathrm{Cat}$). D'autre part, on peut aussi le regarder comme un axiome de \textit{non-trivialité} pour $\mathrm{Hot}_W$. Je rappelle à ce sujet :

\textbf{Proposition 2.7.} \textit{Soit $W$ un localiseur fondamental. Conditions équivalentes :}
\begin{itemize}[leftmargin=3em]
\item[a)] $W$ \textit{satisfait l'axiome de connexité} Loc (5).
\item[b)] \textit{Le foncteur} $\pi_0 : \mathrm{Cat} \longrightarrow \mathrm{Ens}$ \textit{se factorise par un foncteur}
% foot: Version 18 mars 2019
\printedpage{7}
% head: A. Grothendieck : Dérivateurs, Ch. XVI, Édition des Universités de Montpellier II et Paris VII
\[
\begin{tikzcd}
\mathrm{Cat} \arrow[r, dotted] \arrow[rr, dotted, bend right, "\pi_0"'] & \mathrm{Hot}_W \arrow[r, "\pi_0"] & \mathrm{Ens}.
\end{tikzcd}
\]
\item[c)] \textit{La catégorie} $\mathrm{Hot}^*_W = \mathrm{Hot}_W\backslash\{\emptyset\}$ \textit{n'est pas équivalente à la catégorie ponctuelle.}
\item[c$'$)] \textit{La catégorie} $\mathrm{Hot}_W$ \textit{n'est pas équivalente à $e$ ni à $\boldsymbol{\Delta}^1$.}\footnote{\textbf{NB} Si $W$ est fortement saturée (condition Loc (1 ter), alors les conditions ci-contre (\textit{i.e.} Loc (5)) équivalent à ceci : % [sic] "(condition Loc (1 ter)," -> "(condition Loc (1 ter)),", parenthèse ouvrante non refermée
\begin{itemize}[leftmargin=2.5em]
\item[d)] \textit{$W$ n'est égal ni à Fl}(Cat) \textit{tout entier, ni à Fl}(Cat$^*$) $\cup$ $\{\mathrm{id}_\emptyset\}$ (i.e. \textit{Fl}(Cat)$\backslash W$ \textit{n'est ni vide, ni égal à l'ensemble des flèches $\emptyset \longrightarrow X$ avec $X \neq \emptyset$}).
\end{itemize}
Si $W$ satisfait Loc (7 D$_{W_\omega}$), alors $W$ satisfait Loc (5) si et seulement si $W \neq \mathrm{Fl}(\mathrm{Cat})$ (\textit{cf.} XII, th. 2, p. 20).}
\end{itemize}

Pour la démonstration, et pour d'autres formes équivalentes de l'axiome de connexité, \msp{page 13}je renvoie à XV 9.12 (p. 133 bis). Pour le nom “axiome de connexité”, je peux offrir l'explication suivante 2.8 et 2.9.

\textbf{Corollaire 2.8.} \textit{Si l'axiome de connexité} Loc (5) \textit{est satisfait, alors les objets $W$-asphériques sont} 0\textit{-connexes} (\textit{tautologie}), \textit{et la réciproque est vraie si $W$ est fortement saturé.} (\textit{Ou aussi, si $W$ satisfait l'axiome des fibrations} Loc (7) \textit{ci-dessous, puisque par} XII, th. 2, p. 20, \textit{dans ce cas,} Loc (5) \textit{équivaut à $W \neq \mathrm{Fl}(\mathcal{M})$ \ldots}).

Car si l'axiome de connexité n'est pas satisfait, alors pour toute $X$ non vide dans Cat, $X \longrightarrow e$ est un isomorphisme, donc ($W$ étant fortement saturé) $X \longrightarrow e$ est dans $W$, \textit{i.e.} $X$ est $W$-asphérique, donc il y a des objets $W$-asphériques non 0-connexe. Notons aussi : % [sic] "non 0-connexe" -> "non 0-connexes", accord avec "des objets"

\textbf{Corollaire 2.9.} \textit{Supposons que $W$ satisfait à l'axiome} Loc (4) \textit{des sommes directes} (\textit{ou seulement à} Loc (4 b)). \textit{Pour que $W$ satisfasse l'axiome de connexité, il faut et il suffit que pour toute flèche $f : X \longrightarrow Y$ avec $Y$} 0\textit{-connexe, si $f$ est dans $W$, alors $X$ est} 0\textit{-connexe.} (\textit{C'est nécessaire bien sûr sans supposer} Loc (4 b) \textit{ni rien du tout.})

\msp{page 14}\texttt{[Les numéros 2.7, 2.8 et 2.9 apparaissent deux fois.]}

\textbf{2.7.} Considérons maintenant le foncteur composé
\[
\mathrm{Ens} \xrightarrow{I \longmapsto I_{\mathrm{Cat}}} \mathrm{Cat} \xrightarrow{\quad \mathrm{hot}_W \quad} \mathrm{Hot}_W, \qquad I \longmapsto I_{\mathrm{Hot}_W} \text{ ou } I_W, \tag*{(2.7.1)}
\]
où le premier foncteur est celui qui associe à l'ensemble $I$ la catégorie discrète $I_{\mathrm{Cat}}$ ayant $I$ comme ensemble d'objets. (Le plus souvent, on identifiera $I_{\mathrm{Cat}}$ avec $I$ dans les notations, tout comme pour un ensemble ordonné $J$, on identifie $J$ et la “catégorie ordonné” % [sic] "catégorie ordonné" -> "catégorie ordonnée", accord
définie par $J$. Ici $I$ peut être considéré comme muni de l'ordre discret.) On a d'ailleurs
\[
I_{\mathrm{Cat}} \simeq I \times e \ \ (\overset{\text{déf}}{=} \text{somme de $I$ copies de $e$}), \tag*{(2.7.2)}
\]
donc quand Loc (4 a) est satisfait, donc $\mathrm{hot}_W$ commute aux sommes directes, on aura un isomorphisme canonique
\[
\begin{aligned}
I_W \simeq I \times e_W \quad &\text{(où $e_W$ est l'objet final de $\mathrm{Hot}_W$,} \\
&\text{\textit{i.e.} $e$ considéré comme objet de $\mathrm{Hot}_W$).}
\end{aligned} \tag*{(2.7.3)}
\]
% foot: Version 18 mars 2019

\printedpage{8}
% head: A. Grothendieck : Dérivateurs, Ch. XVI, Édition des Universités de Montpellier II et Paris VII
\textbf{Proposition 2.8.} \textit{Supposons l'axiome de connexité} Loc (5) \textit{satisfait, d'où un foncteur}
\[
\pi_0 : \mathrm{Hot}_W \longrightarrow \mathrm{Ens}. \tag*{(2.8.1)}
\]
\textit{Ce foncteur est} adjoint à gauche \textit{du foncteur $I \longmapsto I_W$} (2.7.1), i.e. \textit{on a une bijection fonctorielle en $(X, I)$ dans} $\mathrm{Hot}_W \times \mathrm{Ens}$ :
\msp{page 15}
\[
\mathrm{Hom}_{\mathrm{Ens}}(\pi_0(X), I) \simeq \mathrm{Hom}_W(X, I_W). \tag*{(2.8.2)}
\]
\textit{Le foncteur $\pi_0$ est également} adjoint à droite \textit{du} même \textit{foncteur $I \longmapsto I_W$,} i.e. \textit{on a un bijection fonctorielle} % [sic] "un bijection" -> "une bijection", accord
\[
\mathrm{Hom}_{\mathrm{Ens}}(I, \pi_0(X)) \simeq \mathrm{Hom}_W(I_W, X). \tag*{(2.8.3)}
\]

Prenant $I = e$ (ensemble ponctuel), on trouve le :

\textbf{Corollaire 2.9.} \textit{On a l'isomorphisme fonctoriel en $X$ dans} $\mathrm{Hot}_W$
\[
\boxed{\mathrm{Hom}_W(e_W, X) \simeq \pi_0(X).} \tag*{(2.8.4)}
\]

\textbf{Commentaire 2.10.} La formule d'adjonction (2.8.2) est d'ailleurs familière, chaque fois qu'on a une catégorie $\mathcal{H}$ (telle Cat, où $\mathrm{Hot}_W$ \textit{etc.}) % [sic] "où" -> "ou", conjonction de coordination
munie d'un foncteur $\pi_0$, $\pi_0(X)$ correspondant à l'intuition d'une décomposition d'un objet en “composantes connexes”. (Je ne vais pas expliciter ici les conditions générales sous lesquelles on définit un tel foncteur.) Par contre, l'adjonction en sens inverse de ces deux mêmes foncteurs est un phénomène tout à fait spécial à des catégories \msp{page 16}du type $\mathrm{Hot}_W$. Si on tient compte de la description de $I_W$ comme $I \times e_W$, la formule d'adjonction inhabituelle (2.8.3) équivaut, bien sûr, à la formule (2.8.4) - elle s'en déduit en “élevant à la puissance $I$”. On peut renverser la vapeur, et dans une catégorie générale $\mathcal{H}$, \textit{définir} un foncteur $\pi_0 : \mathcal{H} \longrightarrow \mathrm{Ens}$ par le formule (2.8.4), % [sic] "par le formule" -> "par la formule", accord
\[
\pi_0(X) \overset{\text{déf}}{=} \mathrm{Hom}(e_{\mathcal{H}}, X) \qquad (= \Gamma_{\mathcal{H}}(X), \text{ avec une notation familière}),
\]
où $e_{\mathcal{H}}$ est un objet final de $\mathcal{H}$. On aura donc, en posant $I_{\mathcal{H}} = I \times e_{\mathcal{H}}$, la formule
\[
\mathrm{Hom}(I_{\mathcal{H}}, X) \xrightarrow{\ \sim\ } \mathrm{Hom}(I, \pi_0(X))
\]
(\textit{i.e.} la formule (2.8.3)), disant que l'on a un couple $(I_{\mathcal{H}}, \pi_0)$ de foncteurs adjoints. Supposons d'autre part que pour tout ensemble $I$, l'application canonique d'adjonction
\[
I \longrightarrow \pi_0(I_{\mathcal{H}}) = \mathrm{Hom}_{\mathcal{H}}(e_{\mathcal{H}}, e_{\mathcal{H}} \times I)
\]
soit \textit{bijective}\footnote{ça a tendance à être ultrafaux si $\mathcal{H}$ est additive. Par contre, vrai si $\mathcal{H}$ est une catégorie à sommes distributives par rapport aux produits (voire universelles), dont l'objet final est \textit{connexe} \ldots}, \textit{i.e.} que $I \longmapsto I_{\mathcal{H}}$ soit \textit{pleinement fidèle.} On définit alors, pour tout $I$, une application canonique
% foot: Version 18 mars 2019

\end{document}
